
\textbf{Mise en situation :} Vous apercevez votre enseignant en train de jouer avec son animal de compagnie, Vincent le vélociraptor vorace, sur le terrain de football du cégep. Il lui lance son ballon de basketball préféré que Vincent va chercher en courant avant de le ramener à son maître qui le relance ensuite. Puis, tout à coup, le drame se produit : le ballon éclate.

\textbf{Questions :} Votre cours à option sur les dinosaures abordait justement la notion que la puissance sonore à la source du son produit par un vélociraptor enragé suite à la perte de son ballon était de $\v{P}\text{ W}$. Si vous entendez un son de $\v{Beta}\text{ dB}$, à l'endroit où vous vous trouvez...

\begin{enumerate}[label=\alph*)]
    \item À quelle distance vous trouvez-vous du dinosaure ?
    \item Quelle puissance votre oreille reçoit-elle si on l'assimile à un détecteur circulaire (plat) de rayon $\v{R}\text{ m}$ ?
\end{enumerate}

Votre enseignant essaye de calmer les cris de son dinosaure en faisant jouer la chanson « Everybody Hertz » sur son téléphone cellulaire. Entre les cris du dinosaure (donc alors qu'il reprend son souffle et n'émet pas de son), vous entendez le niveau sonore de la musique à $\v{Musique}\text{ dB}$.

\begin{enumerate}[label= c)]
    \item Quel est le niveau d'intensité sonore totale en dB lorsque vous entendez à la fois la musique et les cris du dinosaure ?
\end{enumerate}